\documentclass[twoside,11pt]{article}

% The build script retrieves the unmodified official JMLR style file from
% the pinned upstream revision when it is not already present locally.
\usepackage[preprint]{jmlr2e}

% Additional packages used by this manuscript. jmlr2e already loads
% amssymb, natbib, graphicx, epsfig, and hyperref.
\usepackage{amsmath,mathtools,bm}
\usepackage{booktabs,array,tabularx,longtable}
\usepackage{enumitem}
\usepackage{tikz}
\usetikzlibrary{arrows.meta,positioning,fit,calc}
\usepackage{xcolor}
\usepackage{microtype}
\usepackage[nameinlink,noabbrev]{cleveref}

\setlist[itemize]{leftmargin=1.5em,itemsep=1.6pt,topsep=2pt}
\setlist[enumerate]{leftmargin=1.8em,itemsep=1.6pt,topsep=2pt}
\renewcommand{\arraystretch}{1.10}

% jmlr2e defines theorem, proposition, lemma, corollary, definition,
% example, remark, and proof. The manuscript additionally uses assumptions.
\newtheorem{assumption}[theorem]{Assumption}

\crefname{assumption}{assumption}{assumptions}
\Crefname{assumption}{Assumption}{Assumptions}
\crefname{lemma}{lemma}{lemmas}
\Crefname{lemma}{Lemma}{Lemmas}
\crefname{proposition}{proposition}{propositions}
\Crefname{proposition}{Proposition}{Propositions}
\crefname{corollary}{corollary}{corollaries}
\Crefname{corollary}{Corollary}{Corollaries}
\crefname{theorem}{theorem}{theorems}
\Crefname{theorem}{Theorem}{Theorems}
\crefname{definition}{definition}{definitions}
\Crefname{definition}{Definition}{Definitions}
\crefname{example}{example}{examples}
\Crefname{example}{Example}{Examples}
\crefname{remark}{remark}{remarks}
\Crefname{remark}{Remark}{Remarks}

\newcommand{\R}{\mathbb{R}}
\newcommand{\N}{\mathbb{N}}
\newcommand{\E}{\mathbb{E}}
\newcommand{\Id}{\operatorname{Id}}
\newcommand{\im}{\operatorname{im}}
\newcommand{\Agg}{\operatorname{Agg}}
\newcommand{\FFN}{\operatorname{FFN}}
\newcommand{\argmin}{\operatorname*{arg\,min}}
\newcommand{\Dep}{\operatorname{Dep}}
\newcommand{\given}{\,\middle|\,}
\newcommand{\norm}[1]{\left\lVert #1\right\rVert}
\newcommand{\abs}[1]{\left|#1\right|}
\newcommand{\ip}[2]{\left\langle #1,#2\right\rangle}
